\documentclass[10pt,a4paper]{amsart}
\usepackage[margin=19mm]{geometry}
\usepackage{amsmath,amssymb,mathtools}
\usepackage[scr=rsfs,cal=euler]{mathalfa}
\usepackage{xfrac}
\usepackage[unicode,hidelinks,bookmarks=false]{hyperref}
\newcommand{\ie}{i.e.\ }
\newcommand{\cf}{cf.\ }
\newcommand{\resp}{respectively\ }
\newcommand{\Subsection}{\subsection}
\providecommand{\nobreakdash}{\nobreak\hbox{-}}
\setlength{\emergencystretch}{2em}
\multlinegap0pt
\usepackage{amssymb}
\usepackage{enumerate}
\usepackage{xcolor}
\newtheorem{theorem}{Theorem}
\newcommand{\fs}{\mathop{\mathrm{FS}}}
\title{full title}
\author{J. A. Corona-Garc\'{\i}a, D. J. Fern\'{a}ndez-Bret\'{o}n, U. A. Ramos-Garc\'{\i}a}
\date{Source: arXiv:2606.09761v1 (CC BY 4.0)}
\begin{document}
\maketitle

\noindent\emph{Source and licence.} full title. Version arXiv:2606.09761v1 (CC BY 4.0), distributed by arXiv under the Creative Commons Attribution 4.0 International licence, http://creativecommons.org/licenses/by/4.0/, as recorded in the arXiv licence metadata for this exact version and shown on the abstract page https://arxiv.org/abs/2606.09761v1. Full licence notice: \texttt{licenses/arxiv-2606.09761-CC-BY-4.0.txt}.\par\smallskip

\noindent\textit{Editorial scope.} Counted unit: the article's theorem thm:arbitrarilylong (source0.tex lines 606--608). The article's complete original proof of it is delivered with it (source0.tex lines 610--638, one \texttt{proof} environment of 29 lines). No mathematics is written, repaired or re-derived: the statement and the proof are the article's own lines, in the article's own notation, and no retained line is substituted or re-flowed.  Retained uncounted as the source-local context the proof needs: lines 559--564.  Nothing was omitted from any retained span, so the package carries no omission record.  No retained line contains a citation, so the wrapper carries no bibliography and no bibliography entry.  No retained line contains a figure environment, an image-inclusion command or drawing code, so no image is delivered and the package declares no asset dependency.  Editorial formatting only, in addition: an AMS \texttt{amsart} preamble that restates the article's own declarations and macros used by the retained lines, namely the environments \texttt{theorem} on separate counters, together with the standard packages those lines need.  The retained environments print in this document as Theorem 1 in order of appearance, whereas the article prints the same items with the numbering of its own full document.  The complete native source of the article as distributed in the arXiv e-print for this exact version is bundled unchanged as \texttt{sources/202426/source0.tex}.
\begin{theorem}\label{thm:centralsets}
    Let $A$ be a central set. Then, whenever we have finitely many sequences $\langle a_n^1\mid n<\omega\rangle,\cdots,\langle a_n^k\mid n<\omega\rangle$, there exists a sequence $\langle x_n\mid n<\omega\rangle$ and a block sequence $\langle H_n\mid n<\omega\rangle$ of finite subsets of $\omega$ such that, for every $t\in\{1,\ldots,k\}$ we have
    \begin{equation*}
        \fs\left(x_n+\sum_{i\in H_t}a_t^i\mid n<\omega\right)\subseteq A.
    \end{equation*}
\end{theorem}

\begin{theorem}\label{thm:arbitrarilylong}
Let $A$ be a central set, and let $l<\omega$. Then there exist sequences $\langle d_n\mid n<\omega\rangle,\langle a_n^1\mid n<\omega\rangle,\ldots,\langle a_n^l\mid n<\omega\rangle$ such that, for all $0\leq i<l,0\leq k\leq l,k-1\leq t\leq k+1$, letting $s_l^n=\langle d_n,a_n^1,\ldots,a_n^l\rangle$ and $\vec{x}^{(i,k,t)}=\langle (i,k,t)*s_l^n\mid n<\omega\rangle$, we have $\fs(\vec{x}^{(i,k,t)})\subseteq A$. In particular (since each element of $L(s^n)$ is the $n$-th term of one of the sequences $\vec{x}^{(i,k,t)}$), the set $A$ contains $l$-patterns for arbitrarily large $l$.
\end{theorem}

\begin{proof}
The proof goes by induction on $l$. For $l=0$ we only need to take a sequence $\langle d_n\mid n<\omega\rangle$ such that $\fs(\langle d_n\mid n<\omega\rangle)\subseteq A$, which is simply Hindman's theorem (since $A$ is central, in particular it is an IP-set).

Now suppose the theorem already holds for some $l$ and let $\langle \delta_n\mid n<\omega\rangle,\langle \alpha_n^1\mid n<\omega\rangle,\ldots,\langle \alpha_n^l\mid n<\omega\rangle$ be the sequences as in the statement of the theorem. We now apply the Central Sets Theorem~\ref{thm:centralsets} to the set $A$ along with the sequences
\begin{equation*}
    \langle(i,k,t)*s_l^n+\delta_n\mid n<\omega\rangle
\end{equation*}
and
\begin{equation*}
    \langle(i,k,t)*s_l^n+2\delta_n\mid n<\omega\rangle
\end{equation*}
for each $(i,k,t)$ such that $1\leq i<l$, $0\leq k\leq l$, and $k-1\leq t\leq k+1$; this yields a sequence $\vec{x}=\langle x_n\mid n<\omega\rangle$ of numbers and a block sequence $\langle H_n\mid n<\omega\rangle$ of finite sets satisfying the statement of Theorem~\ref{thm:centralsets}. Define, for each $n<\omega$,
\begin{eqnarray*}
d_n & = & \sum_{i\in H_n}\delta_i, \\
a_n^j & = & \sum_{i\in H_n}\alpha_i^j, \ \ \ \text{ for }1\leq j\leq l, \\
a_n^{l+1} & = & x_n+\sum_{i\in H_n}\delta_i.
\end{eqnarray*}

We now show that, with $\vec{x}^{(i,k,t)}$ defined as in the statement of the theorem for all of the $l+2$ recently obtained sequences, we have $\fs(\vec{x}^{(i,k,t)})\subseteq A$. If $i+k\leq l$ then the desired result follows immediately by the inductive hypothesis. On the other hand, if $i+k=l+1$, then by looking at the terms of $\vec{x}^{(i,k,t)}$ we obtain 
\begin{equation*}
\vec{x}^{(i,k,t)}=\begin{cases}
\langle x_n+\sum_{i\in H_n}\delta_i\rangle,\ \ \ \text{ if }(i,k,t)=(l,1,0), \\
\langle x_n+\sum_{i\in H_n}2\delta_i\rangle,\ \ \ \text{ if }(i,k,t)=(l,1,1), \\
\langle x_n+\sum_{i\in H_n}\left((i,k-1,t)*s_l^n+\delta_i\right)\rangle,\ \ \ \text{ if }i\leq l\text{ and }t\leq k, \\
\langle x_n+\sum_{i\in H_n}\left((i,k-1,t)*s_l^n+2\delta_i\right)\rangle,\ \ \ \text{ if }i\leq l\text{ and }t=k+1.
\end{cases}
\end{equation*}
All of these sequences have their set of finite sums contained within $A$, by the choice of the sequence of $x_n$ utilizing the Central Sets Theorem.
\end{proof}

\end{document}
